\subsection*{Prayer over the Offerings: gifts recognised before they are given}

The prayer asks God to accept the gifts \emph{which we bring from your own bounty}, and then makes a double petition: that these sacred mysteries, by the working power of grace, may both sanctify the offerers in the conduct of the present life and lead them to everlasting joys.

The governing phrase is checkable and worth checking. \incipit{De tua largit\'ate} -- ``from your bounty'' -- is exactly the phrase the priest has just used, twice, over the bread and over the chalice: blessed are you, Lord God of all creation, for through your goodness we have received this bread, this wine, which earth has given and human hands have made. Four occurrences of the phrase stand in the 2002 Missal's two printings of the Order of Mass, and the Week~XVII Prayer over the Offerings picks it up at the moment those gifts are placed on the altar. The oration is therefore not making a new claim; it is repeating, in the collect form, what the preparation rite has just said, and doing so at the point where the assembly might otherwise imagine it is giving something of its own.

That is a real correction to a possible misreading of the day's Gospel. The two buyers in the parable dispose of property they own. The Prayer over the Offerings says the Church cannot: what it hands over was already a gift. The two statements are not in conflict -- the parable is about the cost of the kingdom to a purchaser, the prayer about the origin of what is offered -- but the liturgy has placed them within a few minutes of each other, and hearing both is more truthful than hearing either alone. \emph{Sacrosanctum concilium}~48 asks precisely this of the faithful: that they learn to offer themselves along with the immaculate victim.

\subsection*{Preface and Eucharistic Prayer: a coupled, unresolved choice}

Week~XVII has no proper Preface and no proper Eucharistic Prayer insert, so two decisions are made locally and neither can be reported here as fact. Under the Missal's Ordinary Time rubrics and GIRM~365 a Preface of Sundays in Ordinary Time is used unless the Eucharistic Prayer selected carries its own; the decisions are therefore coupled. GIRM~365(d) binds Eucharistic Prayer~IV to its own inseparable Preface, excluding a Sunday Preface altogether, while 365(b) leaves Prayer~II free to use its own or another that qualifies.

The consequence is modest but real: the Prefaces of Sundays in Ordinary Time differ in what they say, and none is appointed to these readings. Whatever is heard between the Prayer over the Offerings and the Sanctus is a lawful choice, not a proper text, and no argument here depends on which was made.

\subsection*{Communion Antiphon: a closed pair, unresolved}

The Missal prints two antiphons joined by \emph{Vel}, and one is sung. They pull in different directions, and the difference is worth stating because it changes what the communicant is given to say.

\begin{threecoltable}{Alternative}{What it does}{Direct reception at an exact locus}
\rowthree{A. Ps 102:2 (Vulgate; Ps 103 Hebrew)}{The soul is addressed and told to bless and not to forget the recompenses. It is a command against amnesia, placed at the moment of reception.}{Augustine, \emph{Enarrationes in Psalmos} 103, §2, presses the noun: they are called recompenses, not awards, ``for something else was due, and what was not due has been paid.'' The verse is about gratuity, not merit.}
\rowthree{B. Matthew 5:7--8}{Two adjacent beatitudes, mercy and clean-heartedness, each with its promise. It is a moral test, placed at the same moment.}{Augustine, \emph{On the Lord's Sermon on the Mount} I.2.7--8: the merciful are blessed who relieve the miserable, ``for it is paid back to them''; and a pure heart is a \emph{single} heart, since God is seen with the heart and not with the outward eyes.}
\end{threecoltable}

Neither antiphon is presumed here. Both fit the day without having been chosen for it: A closes a circuit the Prayer over the Offerings opened, since a congregation told that its gifts came from God's bounty is told at Communion not to forget the recompenses; B supplies the disposition that the first reading's request and the Gospel's sorting both imply, without naming anyone. One further connection needs careful bounding. In the same work Augustine correlates the beatitudes with the gifts of the Spirit, assigning \emph{counsel} to the merciful and \emph{understanding} to the clean of heart (I.4.11--12) -- the two gifts a reader of the first reading would name for Solomon's request. The correspondence is Augustine's, made about the Beatitudes and not about this antiphon, and is offered as illumination only.

\subsection*{Prayer after Communion: the one purchase the parables only picture}

The prayer states what has just been received -- a divine sacrament, the \emph{perpetual memorial of the Passion of your Son} -- and asks that this gift, given with ineffable love, may profit our salvation.

The phrase \incipit{passi\'onis F\'ilii tui memori\'ale perp\'etuum} occurs once in the 2002 Missal, here. Its content is the Council's: at the Last Supper Christ instituted the eucharistic sacrifice of his Body and Blood to perpetuate the sacrifice of the Cross through the centuries and to entrust to the Church ``a memorial of His death and resurrection'' (\emph{Sacrosanctum concilium}~47). What the oration adds is the attribution of motive -- the gift was given out of a love the prayer refuses to measure -- and the petition that it be efficacious for salvation rather than merely commemorated.

Set beside the day's Gospel this produces a genuine inversion, and it is a synthesis grounded in the texts rather than a claim about how the Missal was composed. In the parables a man finds a treasure and pays everything he has for it. In the Prayer after Communion the assembly has just received, without price, the memorial of a Passion in which someone else paid. Augustine's answer to the pearl's cost -- that the price is ourselves, ``not because it is worth only that much, but because we cannot give more'' -- sits exactly on the seam between the two, and the Church has always read the sacrament as the place where that inadequate price is accepted.

\subsection*{How far the euchology and the readings may be said to meet}

The three orations are shared by Years A, B and C and are not shown by any evidence available here to have been composed with Lectionary no.~109 in view. The correct classification of the relations below is therefore \emph{source-grounded synthesis}, not official correlation.

\begin{threecoltable}{Point of contact}{Euchological text}{Appointed reading}
\rowthree{Ranking of goods}{Collect: use passing goods now so as already to cling to what abides.}{The psalm ranks commandments above gold and topaz; Solomon declines riches; the buyers convert everything into one object.}
\rowthree{Origin of what is offered}{Prayer over the Offerings: the gifts come from God's own bounty.}{Solomon receives what he asks and also what he does not ask; Romans grounds everything in a divine purpose.}
\rowthree{What is finally received}{Prayer after Communion: the perpetual memorial of the Passion, given out of ineffable love.}{Romans ends in \emph{glorified}; the long Gospel form ends in a sorting no purchaser controls.}
\end{threecoltable}

Two contacts should \emph{not} be claimed. The Collect's ruler-and-guide language is not a comment on Solomon's kingship: it is addressed to God about the petitioners, not about a monarch. And the Entrance antiphon's ``house'' is not the field of the parable; they are different images from different books that happen to fall on the same day.
